\section{Méthode}

% --- Modèle de Bergeron ---
Pour la ligne sans pertes, l'impédance caractéristique, la vitesse de propagation et le temps de
propagation valent
\begin{equation}
Z_c = \sqrt{L'/C'}, \qquad v_p = 1/\sqrt{L'C'}, \qquad \tau = d/v_p ,
\end{equation}
soit $Z_c = \Zc~\unit{\ohm}$, $v_p = \vp$~mi/s et $\tau = \tauL~\unit{\micro s}$.
Le modèle de Bergeron remplace chaque extrémité par la résistance $Z_c$ en parallèle avec une
source de courant d'historique. Cette source ne dépend que de l'autre extrémité à l'instant
$t-\tau$ ; elle est donc connue quand on calcule l'instant $t$. Avec $i_1$ entrant dans la ligne
en $k$ et $i_2$ sortant en $m$ vers la charge :
\begin{align}
i_1(t) &= \frac{v_1(t)}{Z_c} + I_k, & I_k &= i_2(t-\tau) - \frac{v_2(t-\tau)}{Z_c}, \label{eq:k}\\
i_2(t) &= I_m - \frac{v_2(t)}{Z_c}, & I_m &= i_1(t-\tau) + \frac{v_1(t-\tau)}{Z_c}. \label{eq:m}
\end{align}

% --- Modèle compagnon de l'inductance ---
L'inductance de charge est traitée par la règle des trapèzes. L'équation
$\mathrm{d}i_L/\mathrm{d}t = v_L/L$ devient un circuit linéaire, le modèle compagnon : une
résistance $R_{eq}$ en parallèle avec une source $I_h$ connue au pas précédent,
\begin{equation}
i_L(t) = \frac{v_L(t)}{R_{eq}} + I_h, \qquad R_{eq} = \frac{2L}{\Delta t}, \qquad
I_h = i_L(t-\Delta t) + \frac{v_L(t-\Delta t)}{R_{eq}} ,
\end{equation}
avec ici $R_{eq} = \SI{5000}{\ohm}$. Côté $m$, $Z_c \parallel I_m$ alimente $R$ en série avec
$R_{eq} \parallel I_h$. La résolution de ce petit circuit « continu » donne directement
\begin{equation}
i_2 = \frac{I_m + R_{eq} I_h/Z_c}{1 + (R+R_{eq})/Z_c}, \qquad v_2 = Z_c\,(I_m - i_2), \qquad
v_L = v_2 - R\,i_2 .
\end{equation}
Côté $k$, la tension est imposée : $v_1 = v_s$ dès que l'interrupteur est fermé, puis $i_1$ par
\eqref{eq:k}.

% --- Interpolation du retard ---
Le rapport $\tau/\Delta t$ n'est pas entier, donc $t-\tau$ tombe entre deux échantillons
stockés. La valeur retardée s'obtient par interpolation linéaire :
\begin{equation}
x(t-\tau) \approx x_j + a\,(x_{j+1}-x_j), \qquad j = \left\lfloor \frac{t-\tau}{\Delta t}
\right\rfloor, \qquad a = \frac{t-\tau}{\Delta t} - j .
\end{equation}
Il ne faut pas interpoler à travers l'instant de fermeture. Si $t-\tau < t_f$, la ligne était
encore au repos et l'historique vaut zéro. Sinon, l'interpolation entre l'échantillon à $t=0$
(interrupteur ouvert) et celui à $t_f$ ferait arriver l'onde en $m$ un pas trop tôt.
